% ----------------------------------------------------------------------
%  1.3.5. Точечные группы; 1.3.6. Пространственные группы;
%  1.3.7. Предельные группы симметрии
% ----------------------------------------------------------------------
\subsection{Точечные группы симметрии}\label{ssec:point-groups}

\begin{opredelenie}[Точечная группа]
Группа операций симметрии, оставляющих неподвижной хотя бы одну точку.
Ею описывают симметрию тел конечных размеров, а в кристаллографии~—
симметрию внешней формы кристалла и его физических свойств. Точечные
группы кристаллов называют также \term{кристаллическими классами}.
\end{opredelenie}

\paragraph{Группы первого рода}
Из одних поворотных осей получаются 11 групп: циклические $C_1$, $C_2$,
$C_3$, $C_4$, $C_6$ (международные символы 1, 2, 3, 4, 6) и сочетания осей
$222$, $32$, $422$, $622$, $23$, $432$.
\begin{itemize}
\item \emph{Группы поворотов} $N$ ($C_n$): единственный элемент~— ось $N$;
группа циклическая порядка $n=N$. Главная ось полярна.
\item \emph{Группы диэдров} $N2$ ($D_n$): главная ось $N$ и $n$
перпендикулярных ей осей 2; порядок $2n$; главная ось неполярна.
\end{itemize}

Примеры (рядом с порядком группы $n$ указано число классов $k$):
\begin{itemize}
\item $222$ ($D_2$): $E$, $C_2^{[100]}$, $C_2^{[010]}$, $C_2^{[001]}$; $n=4$, $k=4$~—
четверная группа Клейна; генераторы $C_2^{[100]}$ и $C_2^{[010]}$.
\item $422$ ($D_4$): ось 4 и четыре оси 2; $n=8$, $k=5$ (классы: $E$; $C_2$ вокруг
главной оси; два поворота на $\pm90^\circ$; две оси 2 вдоль осей координат;
две диагональные оси 2)\ispr{В оригинале для группы $422$ указано
$k=6$; у группы $D_4$ пять классов сопряжённых элементов.}; генераторы~— $C_4$
и одна из осей 2.
\item $23$ ($T$, группа тетраэдра): четыре оси 3 вдоль $\nfam{111}$ и три оси 2
вдоль $\nfam{100}$; $n=12$, $k=4$ ($E$; три $C_2$; четыре поворота на $120^\circ$;
четыре поворота на $240^\circ$)\ispr{В оригинале для группы $23$ указано
$k=6$; у группы $T$ четыре класса: повороты на $120^\circ$ и на $240^\circ$
вокруг осей $\nfam{111}$ образуют два \emph{разных} класса.}; генераторы
$C_2^{[001]}$ и $C_3^{[111]}$.
\item $432$ ($O$, группа октаэдра): три оси 4, четыре оси 3, шесть осей 2;
$n=24$, $k=5$; генераторы $C_4^{[001]}$ и $C_3^{[111]}$.
\end{itemize}

Сравнение $T$ и $O$ поучительно: равенство углов поворота~— необходимое,
но не достаточное условие сопряжённости. Восемь поворотов на $\pm120^\circ$
вокруг $\nfam{111}$ в группе $O$ образуют один класс, а в группе $T$~— два (в
$T$ нет операции, переводящей поворот на $+120^\circ$ в поворот на
$-120^\circ$).

\paragraph{Все 32 класса}
Добавив инверсионные оси, получаем ещё 21 группу (второго рода)~— всего
\vazhno{32 точечные группы} (табл.~\ref{tab:point-groups}). Самая
симметричная, $m\br3m$ ($O_h$),~— группа куба: 48 элементов, 10 классов,
генераторы $C_4^{[001]}$, $C_3^{[111]}$ и $\br1$ (рис.~\ref{fig:m3m}).

\begin{table}[p]
\centering
\caption{Кристаллические классы (точечные группы) по сингониям}\label{tab:point-groups}
\small\renewcommand{\arraystretch}{1.08}
\begin{tabular}{@{}llll@{}}
\toprule
Сингония & Международный & Шёнфлис & Формула симметрии\\
\midrule
Триклинная & $1$ & $C_1$ & $L_1$\\
 & $\br1$ & $C_i=S_2$ & $C$\\
\midrule
Моноклинная & $2$ & $C_2$ & $L_2$\\
 & $m$ & $C_s=C_{1h}$ & $P$\\
 & $2/m$ & $C_{2h}$ & $L_2PC$\\
\midrule
Ромбическая & $222$ & $D_2=V$ & $3L_2$\\
 & $mm2$ & $C_{2v}$ & $L_22P$\\
 & $mmm$ & $D_{2h}=V_h$ & $3L_23PC$\\
\midrule
Тетрагональная & $4$ & $C_4$ & $L_4$\\
 & $\br4$ & $S_4$ & $\bar L_4$\\
 & $4/m$ & $C_{4h}$ & $L_4PC$\\
 & $422$ & $D_4$ & $L_44L_2$\\
 & $4mm$ & $C_{4v}$ & $L_44P$\\
 & $\br42m$ & $D_{2d}=V_d$ & $\bar L_42L_22P$\\
 & $4/mmm$ & $D_{4h}$ & $L_44L_25PC$\\
\midrule
Тригональная & $3$ & $C_3$ & $L_3$\\
 & $\br3$ & $C_{3i}=S_6$ & $L_3C$\\
 & $32$ & $D_3$ & $L_33L_2$\\
 & $3m$ & $C_{3v}$ & $L_33P$\\
 & $\br3m$ & $D_{3d}$ & $L_33L_23PC$\\
\midrule
Гексагональная & $6$ & $C_6$ & $L_6$\\
 & $\br6$ & $C_{3h}$ & $L_3P$\\
 & $6/m$ & $C_{6h}$ & $L_6PC$\\
 & $622$ & $D_6$ & $L_66L_2$\\
 & $6mm$ & $C_{6v}$ & $L_66P$\\
 & $\br6m2$ & $D_{3h}$ & $L_33L_24P$\\
 & $6/mmm$ & $D_{6h}$ & $L_66L_27PC$\\
\midrule
Кубическая & $23$ & $T$ & $4L_33L_2$\\
 & $m\br3$ & $T_h$ & $4L_33L_23PC$\\
 & $432$ & $O$ & $3L_44L_36L_2$\\
 & $\br43m$ & $T_d$ & $3\bar L_44L_36P$\\
 & $m\br3m$ & $O_h$ & $3L_44L_36L_29PC$\\
\bottomrule
\multicolumn{4}{p{0.8\textwidth}}{\footnotesize В формуле симметрии $L_n$~— ось $n$-го порядка, $\bar L_n$~— инверсионная ось, $P$~— плоскость симметрии, $C$~— центр симметрии; число перед символом~— количество элементов.}
\end{tabular}
\end{table}

\paragraph{Обозначения}
В международной символике плоскость, перпендикулярная оси $N$, пишется
через косую черту ($N/m$), а плоскость, содержащая ось,~— рядом ($Nm$).
У Шёнфлиса первым соответствуют группы $C_{nh}$ (h~— горизонтальная
плоскость), вторым~— $C_{nv}$ (v~— вертикальная), $D_n$~— группы диэдров,
$S_n$~— зеркально-поворотные оси, $T$ и $O$~— группы тетраэдра и
октаэдра. В символике Шубникова точка между элементами означает их
параллельность, двоеточие~— перпендикулярность, косая черта~— косое
расположение.

\begin{figure}[ht]
\centering
\tikzset{mirror/.style={fmBlue, line width=0.6pt},
  ax4/.style={rectangle, draw=fmRed, fill=fmRed!35, inner sep=0pt, minimum size=5pt},
  ax3/.style={regular polygon, regular polygon sides=3, draw=fmAmber!90!black, fill=fmAmber!40, inner sep=0pt, minimum size=7pt},
  ax2/.style={ellipse, draw=fmGreen!60!black, fill=fmGreen!35, inner sep=0pt, minimum width=7pt, minimum height=4pt}}
\begin{tikzpicture}[scale=2.4]
  \draw[mirror, line width=1pt] (0,0) circle (1);
  % создано tools/gen_stereo.py: элементы симметрии m-3m на проекции (001)
\draw[mirror] (0.0000,-1.0000) -- (0.0000,-0.9742) -- (0.0000,-0.9490) -- (0.0000,-0.9244) -- (0.0000,-0.9004) -- (0.0000,-0.8770) -- (0.0000,-0.8541) -- (0.0000,-0.8317) -- (0.0000,-0.8098) -- (0.0000,-0.7883) -- (0.0000,-0.7673) -- (0.0000,-0.7467) -- (0.0000,-0.7265) -- (0.0000,-0.7067) -- (0.0000,-0.6873) -- (0.0000,-0.6682) -- (0.0000,-0.6494) -- (0.0000,-0.6310) -- (0.0000,-0.6128) -- (0.0000,-0.5949) -- (0.0000,-0.5774) -- (0.0000,-0.5600) -- (0.0000,-0.5430) -- (0.0000,-0.5261) -- (0.0000,-0.5095) -- (0.0000,-0.4931) -- (0.0000,-0.4770) -- (0.0000,-0.4610) -- (0.0000,-0.4452) -- (0.0000,-0.4296) -- (0.0000,-0.4142) -- (0.0000,-0.3990) -- (0.0000,-0.3839) -- (0.0000,-0.3689) -- (0.0000,-0.3541) -- (0.0000,-0.3395) -- (0.0000,-0.3249) -- (0.0000,-0.3105) -- (0.0000,-0.2962) -- (0.0000,-0.2820) -- (0.0000,-0.2679) -- (0.0000,-0.2540) -- (0.0000,-0.2401) -- (0.0000,-0.2263) -- (0.0000,-0.2126) -- (0.0000,-0.1989) -- (0.0000,-0.1853) -- (0.0000,-0.1718) -- (0.0000,-0.1584) -- (0.0000,-0.1450) -- (0.0000,-0.1317) -- (0.0000,-0.1184) -- (0.0000,-0.1051) -- (0.0000,-0.0919) -- (0.0000,-0.0787) -- (0.0000,-0.0655) -- (0.0000,-0.0524) -- (0.0000,-0.0393) -- (0.0000,-0.0262) -- (0.0000,-0.0131) -- (0.0000,-0.0000) -- (-0.0000,0.0131) -- (-0.0000,0.0262) -- (-0.0000,0.0393) -- (-0.0000,0.0524) -- (-0.0000,0.0655) -- (-0.0000,0.0787) -- (-0.0000,0.0919) -- (-0.0000,0.1051) -- (-0.0000,0.1184) -- (-0.0000,0.1317) -- (-0.0000,0.1450) -- (-0.0000,0.1584) -- (-0.0000,0.1718) -- (-0.0000,0.1853) -- (-0.0000,0.1989) -- (-0.0000,0.2126) -- (-0.0000,0.2263) -- (-0.0000,0.2401) -- (-0.0000,0.2540) -- (-0.0000,0.2679) -- (-0.0000,0.2820) -- (-0.0000,0.2962) -- (-0.0000,0.3105) -- (-0.0000,0.3249) -- (-0.0000,0.3395) -- (-0.0000,0.3541) -- (-0.0000,0.3689) -- (-0.0000,0.3839) -- (-0.0000,0.3990) -- (-0.0000,0.4142) -- (-0.0000,0.4296) -- (-0.0000,0.4452) -- (-0.0000,0.4610) -- (-0.0000,0.4770) -- (-0.0000,0.4931) -- (-0.0000,0.5095) -- (-0.0000,0.5261) -- (-0.0000,0.5430) -- (-0.0000,0.5600) -- (-0.0000,0.5774) -- (-0.0000,0.5949) -- (-0.0000,0.6128) -- (-0.0000,0.6310) -- (-0.0000,0.6494) -- (-0.0000,0.6682) -- (-0.0000,0.6873) -- (-0.0000,0.7067) -- (-0.0000,0.7265) -- (-0.0000,0.7467) -- (-0.0000,0.7673) -- (-0.0000,0.7883) -- (-0.0000,0.8098) -- (-0.0000,0.8317) -- (-0.0000,0.8541) -- (-0.0000,0.8770) -- (-0.0000,0.9004) -- (-0.0000,0.9244) -- (-0.0000,0.9490) -- (-0.0000,0.9742) -- (-0.0000,1.0000);
\draw[mirror] (1.0000,0.0000) -- (0.9742,0.0000) -- (0.9490,0.0000) -- (0.9244,0.0000) -- (0.9004,0.0000) -- (0.8770,0.0000) -- (0.8541,0.0000) -- (0.8317,0.0000) -- (0.8098,0.0000) -- (0.7883,0.0000) -- (0.7673,0.0000) -- (0.7467,0.0000) -- (0.7265,0.0000) -- (0.7067,0.0000) -- (0.6873,0.0000) -- (0.6682,0.0000) -- (0.6494,0.0000) -- (0.6310,0.0000) -- (0.6128,0.0000) -- (0.5949,0.0000) -- (0.5774,0.0000) -- (0.5600,0.0000) -- (0.5430,0.0000) -- (0.5261,0.0000) -- (0.5095,0.0000) -- (0.4931,0.0000) -- (0.4770,0.0000) -- (0.4610,0.0000) -- (0.4452,0.0000) -- (0.4296,0.0000) -- (0.4142,0.0000) -- (0.3990,0.0000) -- (0.3839,0.0000) -- (0.3689,0.0000) -- (0.3541,0.0000) -- (0.3395,0.0000) -- (0.3249,0.0000) -- (0.3105,0.0000) -- (0.2962,0.0000) -- (0.2820,0.0000) -- (0.2679,0.0000) -- (0.2540,0.0000) -- (0.2401,0.0000) -- (0.2263,0.0000) -- (0.2126,0.0000) -- (0.1989,0.0000) -- (0.1853,0.0000) -- (0.1718,0.0000) -- (0.1584,0.0000) -- (0.1450,0.0000) -- (0.1317,0.0000) -- (0.1184,0.0000) -- (0.1051,0.0000) -- (0.0919,0.0000) -- (0.0787,0.0000) -- (0.0655,0.0000) -- (0.0524,0.0000) -- (0.0393,0.0000) -- (0.0262,0.0000) -- (0.0131,0.0000) -- (0.0000,0.0000) -- (-0.0131,-0.0000) -- (-0.0262,-0.0000) -- (-0.0393,-0.0000) -- (-0.0524,-0.0000) -- (-0.0655,-0.0000) -- (-0.0787,-0.0000) -- (-0.0919,-0.0000) -- (-0.1051,-0.0000) -- (-0.1184,-0.0000) -- (-0.1317,-0.0000) -- (-0.1450,-0.0000) -- (-0.1584,-0.0000) -- (-0.1718,-0.0000) -- (-0.1853,-0.0000) -- (-0.1989,-0.0000) -- (-0.2126,-0.0000) -- (-0.2263,-0.0000) -- (-0.2401,-0.0000) -- (-0.2540,-0.0000) -- (-0.2679,-0.0000) -- (-0.2820,-0.0000) -- (-0.2962,-0.0000) -- (-0.3105,-0.0000) -- (-0.3249,-0.0000) -- (-0.3395,-0.0000) -- (-0.3541,-0.0000) -- (-0.3689,-0.0000) -- (-0.3839,-0.0000) -- (-0.3990,-0.0000) -- (-0.4142,-0.0000) -- (-0.4296,-0.0000) -- (-0.4452,-0.0000) -- (-0.4610,-0.0000) -- (-0.4770,-0.0000) -- (-0.4931,-0.0000) -- (-0.5095,-0.0000) -- (-0.5261,-0.0000) -- (-0.5430,-0.0000) -- (-0.5600,-0.0000) -- (-0.5774,-0.0000) -- (-0.5949,-0.0000) -- (-0.6128,-0.0000) -- (-0.6310,-0.0000) -- (-0.6494,-0.0000) -- (-0.6682,-0.0000) -- (-0.6873,-0.0000) -- (-0.7067,-0.0000) -- (-0.7265,-0.0000) -- (-0.7467,-0.0000) -- (-0.7673,-0.0000) -- (-0.7883,-0.0000) -- (-0.8098,-0.0000) -- (-0.8317,-0.0000) -- (-0.8541,-0.0000) -- (-0.8770,-0.0000) -- (-0.9004,-0.0000) -- (-0.9244,-0.0000) -- (-0.9490,-0.0000) -- (-0.9742,-0.0000) -- (-1.0000,-0.0000);
\draw[mirror] (0.7071,-0.7071) -- (0.6888,-0.6888) -- (0.6710,-0.6710) -- (0.6536,-0.6536) -- (0.6367,-0.6367) -- (0.6201,-0.6201) -- (0.6039,-0.6039) -- (0.5881,-0.5881) -- (0.5726,-0.5726) -- (0.5574,-0.5574) -- (0.5426,-0.5426) -- (0.5280,-0.5280) -- (0.5137,-0.5137) -- (0.4997,-0.4997) -- (0.4860,-0.4860) -- (0.4725,-0.4725) -- (0.4592,-0.4592) -- (0.4462,-0.4462) -- (0.4333,-0.4333) -- (0.4207,-0.4207) -- (0.4082,-0.4082) -- (0.3960,-0.3960) -- (0.3839,-0.3839) -- (0.3720,-0.3720) -- (0.3603,-0.3603) -- (0.3487,-0.3487) -- (0.3373,-0.3373) -- (0.3260,-0.3260) -- (0.3148,-0.3148) -- (0.3038,-0.3038) -- (0.2929,-0.2929) -- (0.2821,-0.2821) -- (0.2714,-0.2714) -- (0.2609,-0.2609) -- (0.2504,-0.2504) -- (0.2400,-0.2400) -- (0.2298,-0.2298) -- (0.2196,-0.2196) -- (0.2095,-0.2095) -- (0.1994,-0.1994) -- (0.1895,-0.1895) -- (0.1796,-0.1796) -- (0.1698,-0.1698) -- (0.1600,-0.1600) -- (0.1503,-0.1503) -- (0.1407,-0.1407) -- (0.1311,-0.1311) -- (0.1215,-0.1215) -- (0.1120,-0.1120) -- (0.1025,-0.1025) -- (0.0931,-0.0931) -- (0.0837,-0.0837) -- (0.0743,-0.0743) -- (0.0650,-0.0650) -- (0.0557,-0.0557) -- (0.0463,-0.0463) -- (0.0371,-0.0371) -- (0.0278,-0.0278) -- (0.0185,-0.0185) -- (0.0093,-0.0093) -- (0.0000,-0.0000) -- (-0.0093,0.0093) -- (-0.0185,0.0185) -- (-0.0278,0.0278) -- (-0.0371,0.0371) -- (-0.0463,0.0463) -- (-0.0557,0.0557) -- (-0.0650,0.0650) -- (-0.0743,0.0743) -- (-0.0837,0.0837) -- (-0.0931,0.0931) -- (-0.1025,0.1025) -- (-0.1120,0.1120) -- (-0.1215,0.1215) -- (-0.1311,0.1311) -- (-0.1407,0.1407) -- (-0.1503,0.1503) -- (-0.1600,0.1600) -- (-0.1698,0.1698) -- (-0.1796,0.1796) -- (-0.1895,0.1895) -- (-0.1994,0.1994) -- (-0.2095,0.2095) -- (-0.2196,0.2196) -- (-0.2298,0.2298) -- (-0.2400,0.2400) -- (-0.2504,0.2504) -- (-0.2609,0.2609) -- (-0.2714,0.2714) -- (-0.2821,0.2821) -- (-0.2929,0.2929) -- (-0.3038,0.3038) -- (-0.3148,0.3148) -- (-0.3260,0.3260) -- (-0.3373,0.3373) -- (-0.3487,0.3487) -- (-0.3603,0.3603) -- (-0.3720,0.3720) -- (-0.3839,0.3839) -- (-0.3960,0.3960) -- (-0.4082,0.4082) -- (-0.4207,0.4207) -- (-0.4333,0.4333) -- (-0.4462,0.4462) -- (-0.4592,0.4592) -- (-0.4725,0.4725) -- (-0.4860,0.4860) -- (-0.4997,0.4997) -- (-0.5137,0.5137) -- (-0.5280,0.5280) -- (-0.5426,0.5426) -- (-0.5574,0.5574) -- (-0.5726,0.5726) -- (-0.5881,0.5881) -- (-0.6039,0.6039) -- (-0.6201,0.6201) -- (-0.6367,0.6367) -- (-0.6536,0.6536) -- (-0.6710,0.6710) -- (-0.6888,0.6888) -- (-0.7071,0.7071);
\draw[mirror] (-0.7071,-0.7071) -- (-0.6888,-0.6888) -- (-0.6710,-0.6710) -- (-0.6536,-0.6536) -- (-0.6367,-0.6367) -- (-0.6201,-0.6201) -- (-0.6039,-0.6039) -- (-0.5881,-0.5881) -- (-0.5726,-0.5726) -- (-0.5574,-0.5574) -- (-0.5426,-0.5426) -- (-0.5280,-0.5280) -- (-0.5137,-0.5137) -- (-0.4997,-0.4997) -- (-0.4860,-0.4860) -- (-0.4725,-0.4725) -- (-0.4592,-0.4592) -- (-0.4462,-0.4462) -- (-0.4333,-0.4333) -- (-0.4207,-0.4207) -- (-0.4082,-0.4082) -- (-0.3960,-0.3960) -- (-0.3839,-0.3839) -- (-0.3720,-0.3720) -- (-0.3603,-0.3603) -- (-0.3487,-0.3487) -- (-0.3373,-0.3373) -- (-0.3260,-0.3260) -- (-0.3148,-0.3148) -- (-0.3038,-0.3038) -- (-0.2929,-0.2929) -- (-0.2821,-0.2821) -- (-0.2714,-0.2714) -- (-0.2609,-0.2609) -- (-0.2504,-0.2504) -- (-0.2400,-0.2400) -- (-0.2298,-0.2298) -- (-0.2196,-0.2196) -- (-0.2095,-0.2095) -- (-0.1994,-0.1994) -- (-0.1895,-0.1895) -- (-0.1796,-0.1796) -- (-0.1698,-0.1698) -- (-0.1600,-0.1600) -- (-0.1503,-0.1503) -- (-0.1407,-0.1407) -- (-0.1311,-0.1311) -- (-0.1215,-0.1215) -- (-0.1120,-0.1120) -- (-0.1025,-0.1025) -- (-0.0931,-0.0931) -- (-0.0837,-0.0837) -- (-0.0743,-0.0743) -- (-0.0650,-0.0650) -- (-0.0557,-0.0557) -- (-0.0463,-0.0463) -- (-0.0371,-0.0371) -- (-0.0278,-0.0278) -- (-0.0185,-0.0185) -- (-0.0093,-0.0093) -- (-0.0000,-0.0000) -- (0.0093,0.0093) -- (0.0185,0.0185) -- (0.0278,0.0278) -- (0.0371,0.0371) -- (0.0463,0.0463) -- (0.0557,0.0557) -- (0.0650,0.0650) -- (0.0743,0.0743) -- (0.0837,0.0837) -- (0.0931,0.0931) -- (0.1025,0.1025) -- (0.1120,0.1120) -- (0.1215,0.1215) -- (0.1311,0.1311) -- (0.1407,0.1407) -- (0.1503,0.1503) -- (0.1600,0.1600) -- (0.1698,0.1698) -- (0.1796,0.1796) -- (0.1895,0.1895) -- (0.1994,0.1994) -- (0.2095,0.2095) -- (0.2196,0.2196) -- (0.2298,0.2298) -- (0.2400,0.2400) -- (0.2504,0.2504) -- (0.2609,0.2609) -- (0.2714,0.2714) -- (0.2821,0.2821) -- (0.2929,0.2929) -- (0.3038,0.3038) -- (0.3148,0.3148) -- (0.3260,0.3260) -- (0.3373,0.3373) -- (0.3487,0.3487) -- (0.3603,0.3603) -- (0.3720,0.3720) -- (0.3839,0.3839) -- (0.3960,0.3960) -- (0.4082,0.4082) -- (0.4207,0.4207) -- (0.4333,0.4333) -- (0.4462,0.4462) -- (0.4592,0.4592) -- (0.4725,0.4725) -- (0.4860,0.4860) -- (0.4997,0.4997) -- (0.5137,0.5137) -- (0.5280,0.5280) -- (0.5426,0.5426) -- (0.5574,0.5574) -- (0.5726,0.5726) -- (0.5881,0.5881) -- (0.6039,0.6039) -- (0.6201,0.6201) -- (0.6367,0.6367) -- (0.6536,0.6536) -- (0.6710,0.6710) -- (0.6888,0.6888) -- (0.7071,0.7071);
\draw[mirror] (0.0000,-1.0000) -- (-0.0182,-0.9815) -- (-0.0357,-0.9630) -- (-0.0526,-0.9445) -- (-0.0688,-0.9261) -- (-0.0845,-0.9077) -- (-0.0996,-0.8893) -- (-0.1142,-0.8710) -- (-0.1282,-0.8528) -- (-0.1417,-0.8346) -- (-0.1547,-0.8165) -- (-0.1672,-0.7985) -- (-0.1793,-0.7805) -- (-0.1910,-0.7626) -- (-0.2022,-0.7448) -- (-0.2130,-0.7271) -- (-0.2234,-0.7095) -- (-0.2334,-0.6919) -- (-0.2430,-0.6745) -- (-0.2523,-0.6571) -- (-0.2612,-0.6398) -- (-0.2698,-0.6226) -- (-0.2780,-0.6055) -- (-0.2860,-0.5884) -- (-0.2936,-0.5715) -- (-0.3009,-0.5546) -- (-0.3080,-0.5378) -- (-0.3147,-0.5211) -- (-0.3212,-0.5045) -- (-0.3274,-0.4879) -- (-0.3333,-0.4714) -- (-0.3390,-0.4550) -- (-0.3445,-0.4386) -- (-0.3497,-0.4224) -- (-0.3546,-0.4061) -- (-0.3594,-0.3900) -- (-0.3639,-0.3739) -- (-0.3682,-0.3579) -- (-0.3723,-0.3419) -- (-0.3761,-0.3260) -- (-0.3798,-0.3101) -- (-0.3833,-0.2943) -- (-0.3865,-0.2785) -- (-0.3896,-0.2628) -- (-0.3925,-0.2471) -- (-0.3951,-0.2315) -- (-0.3976,-0.2159) -- (-0.4000,-0.2003) -- (-0.4021,-0.1848) -- (-0.4040,-0.1693) -- (-0.4058,-0.1538) -- (-0.4074,-0.1383) -- (-0.4089,-0.1229) -- (-0.4101,-0.1075) -- (-0.4112,-0.0921) -- (-0.4121,-0.0767) -- (-0.4129,-0.0614) -- (-0.4135,-0.0460) -- (-0.4139,-0.0307) -- (-0.4141,-0.0153) -- (-0.4142,-0.0000) -- (-0.4141,0.0153) -- (-0.4139,0.0307) -- (-0.4135,0.0460) -- (-0.4129,0.0614) -- (-0.4121,0.0767) -- (-0.4112,0.0921) -- (-0.4101,0.1075) -- (-0.4089,0.1229) -- (-0.4074,0.1383) -- (-0.4058,0.1538) -- (-0.4040,0.1693) -- (-0.4021,0.1848) -- (-0.4000,0.2003) -- (-0.3976,0.2159) -- (-0.3951,0.2315) -- (-0.3925,0.2471) -- (-0.3896,0.2628) -- (-0.3865,0.2785) -- (-0.3833,0.2943) -- (-0.3798,0.3101) -- (-0.3761,0.3260) -- (-0.3723,0.3419) -- (-0.3682,0.3579) -- (-0.3639,0.3739) -- (-0.3594,0.3900) -- (-0.3546,0.4061) -- (-0.3497,0.4224) -- (-0.3445,0.4386) -- (-0.3390,0.4550) -- (-0.3333,0.4714) -- (-0.3274,0.4879) -- (-0.3212,0.5045) -- (-0.3147,0.5211) -- (-0.3080,0.5378) -- (-0.3009,0.5546) -- (-0.2936,0.5715) -- (-0.2860,0.5884) -- (-0.2780,0.6055) -- (-0.2698,0.6226) -- (-0.2612,0.6398) -- (-0.2523,0.6571) -- (-0.2430,0.6745) -- (-0.2334,0.6919) -- (-0.2234,0.7095) -- (-0.2130,0.7271) -- (-0.2022,0.7448) -- (-0.1910,0.7626) -- (-0.1793,0.7805) -- (-0.1672,0.7985) -- (-0.1547,0.8165) -- (-0.1417,0.8346) -- (-0.1282,0.8528) -- (-0.1142,0.8710) -- (-0.0996,0.8893) -- (-0.0845,0.9077) -- (-0.0688,0.9261) -- (-0.0526,0.9445) -- (-0.0357,0.9630) -- (-0.0182,0.9815) -- (-0.0000,1.0000);
\draw[mirror] (0.0000,-1.0000) -- (0.0182,-0.9815) -- (0.0357,-0.9630) -- (0.0526,-0.9445) -- (0.0688,-0.9261) -- (0.0845,-0.9077) -- (0.0996,-0.8893) -- (0.1142,-0.8710) -- (0.1282,-0.8528) -- (0.1417,-0.8346) -- (0.1547,-0.8165) -- (0.1672,-0.7985) -- (0.1793,-0.7805) -- (0.1910,-0.7626) -- (0.2022,-0.7448) -- (0.2130,-0.7271) -- (0.2234,-0.7095) -- (0.2334,-0.6919) -- (0.2430,-0.6745) -- (0.2523,-0.6571) -- (0.2612,-0.6398) -- (0.2698,-0.6226) -- (0.2780,-0.6055) -- (0.2860,-0.5884) -- (0.2936,-0.5715) -- (0.3009,-0.5546) -- (0.3080,-0.5378) -- (0.3147,-0.5211) -- (0.3212,-0.5045) -- (0.3274,-0.4879) -- (0.3333,-0.4714) -- (0.3390,-0.4550) -- (0.3445,-0.4386) -- (0.3497,-0.4224) -- (0.3546,-0.4061) -- (0.3594,-0.3900) -- (0.3639,-0.3739) -- (0.3682,-0.3579) -- (0.3723,-0.3419) -- (0.3761,-0.3260) -- (0.3798,-0.3101) -- (0.3833,-0.2943) -- (0.3865,-0.2785) -- (0.3896,-0.2628) -- (0.3925,-0.2471) -- (0.3951,-0.2315) -- (0.3976,-0.2159) -- (0.4000,-0.2003) -- (0.4021,-0.1848) -- (0.4040,-0.1693) -- (0.4058,-0.1538) -- (0.4074,-0.1383) -- (0.4089,-0.1229) -- (0.4101,-0.1075) -- (0.4112,-0.0921) -- (0.4121,-0.0767) -- (0.4129,-0.0614) -- (0.4135,-0.0460) -- (0.4139,-0.0307) -- (0.4141,-0.0153) -- (0.4142,-0.0000) -- (0.4141,0.0153) -- (0.4139,0.0307) -- (0.4135,0.0460) -- (0.4129,0.0614) -- (0.4121,0.0767) -- (0.4112,0.0921) -- (0.4101,0.1075) -- (0.4089,0.1229) -- (0.4074,0.1383) -- (0.4058,0.1538) -- (0.4040,0.1693) -- (0.4021,0.1848) -- (0.4000,0.2003) -- (0.3976,0.2159) -- (0.3951,0.2315) -- (0.3925,0.2471) -- (0.3896,0.2628) -- (0.3865,0.2785) -- (0.3833,0.2943) -- (0.3798,0.3101) -- (0.3761,0.3260) -- (0.3723,0.3419) -- (0.3682,0.3579) -- (0.3639,0.3739) -- (0.3594,0.3900) -- (0.3546,0.4061) -- (0.3497,0.4224) -- (0.3445,0.4386) -- (0.3390,0.4550) -- (0.3333,0.4714) -- (0.3274,0.4879) -- (0.3212,0.5045) -- (0.3147,0.5211) -- (0.3080,0.5378) -- (0.3009,0.5546) -- (0.2936,0.5715) -- (0.2860,0.5884) -- (0.2780,0.6055) -- (0.2698,0.6226) -- (0.2612,0.6398) -- (0.2523,0.6571) -- (0.2430,0.6745) -- (0.2334,0.6919) -- (0.2234,0.7095) -- (0.2130,0.7271) -- (0.2022,0.7448) -- (0.1910,0.7626) -- (0.1793,0.7805) -- (0.1672,0.7985) -- (0.1547,0.8165) -- (0.1417,0.8346) -- (0.1282,0.8528) -- (0.1142,0.8710) -- (0.0996,0.8893) -- (0.0845,0.9077) -- (0.0688,0.9261) -- (0.0526,0.9445) -- (0.0357,0.9630) -- (0.0182,0.9815) -- (0.0000,1.0000);
\draw[mirror] (1.0000,0.0000) -- (0.9815,-0.0182) -- (0.9630,-0.0357) -- (0.9445,-0.0526) -- (0.9261,-0.0688) -- (0.9077,-0.0845) -- (0.8893,-0.0996) -- (0.8710,-0.1142) -- (0.8528,-0.1282) -- (0.8346,-0.1417) -- (0.8165,-0.1547) -- (0.7985,-0.1672) -- (0.7805,-0.1793) -- (0.7626,-0.1910) -- (0.7448,-0.2022) -- (0.7271,-0.2130) -- (0.7095,-0.2234) -- (0.6919,-0.2334) -- (0.6745,-0.2430) -- (0.6571,-0.2523) -- (0.6398,-0.2612) -- (0.6226,-0.2698) -- (0.6055,-0.2780) -- (0.5884,-0.2860) -- (0.5715,-0.2936) -- (0.5546,-0.3009) -- (0.5378,-0.3080) -- (0.5211,-0.3147) -- (0.5045,-0.3212) -- (0.4879,-0.3274) -- (0.4714,-0.3333) -- (0.4550,-0.3390) -- (0.4386,-0.3445) -- (0.4224,-0.3497) -- (0.4061,-0.3546) -- (0.3900,-0.3594) -- (0.3739,-0.3639) -- (0.3579,-0.3682) -- (0.3419,-0.3723) -- (0.3260,-0.3761) -- (0.3101,-0.3798) -- (0.2943,-0.3833) -- (0.2785,-0.3865) -- (0.2628,-0.3896) -- (0.2471,-0.3925) -- (0.2315,-0.3951) -- (0.2159,-0.3976) -- (0.2003,-0.4000) -- (0.1848,-0.4021) -- (0.1693,-0.4040) -- (0.1538,-0.4058) -- (0.1383,-0.4074) -- (0.1229,-0.4089) -- (0.1075,-0.4101) -- (0.0921,-0.4112) -- (0.0767,-0.4121) -- (0.0614,-0.4129) -- (0.0460,-0.4135) -- (0.0307,-0.4139) -- (0.0153,-0.4141) -- (0.0000,-0.4142) -- (-0.0153,-0.4141) -- (-0.0307,-0.4139) -- (-0.0460,-0.4135) -- (-0.0614,-0.4129) -- (-0.0767,-0.4121) -- (-0.0921,-0.4112) -- (-0.1075,-0.4101) -- (-0.1229,-0.4089) -- (-0.1383,-0.4074) -- (-0.1538,-0.4058) -- (-0.1693,-0.4040) -- (-0.1848,-0.4021) -- (-0.2003,-0.4000) -- (-0.2159,-0.3976) -- (-0.2315,-0.3951) -- (-0.2471,-0.3925) -- (-0.2628,-0.3896) -- (-0.2785,-0.3865) -- (-0.2943,-0.3833) -- (-0.3101,-0.3798) -- (-0.3260,-0.3761) -- (-0.3419,-0.3723) -- (-0.3579,-0.3682) -- (-0.3739,-0.3639) -- (-0.3900,-0.3594) -- (-0.4061,-0.3546) -- (-0.4224,-0.3497) -- (-0.4386,-0.3445) -- (-0.4550,-0.3390) -- (-0.4714,-0.3333) -- (-0.4879,-0.3274) -- (-0.5045,-0.3212) -- (-0.5211,-0.3147) -- (-0.5378,-0.3080) -- (-0.5546,-0.3009) -- (-0.5715,-0.2936) -- (-0.5884,-0.2860) -- (-0.6055,-0.2780) -- (-0.6226,-0.2698) -- (-0.6398,-0.2612) -- (-0.6571,-0.2523) -- (-0.6745,-0.2430) -- (-0.6919,-0.2334) -- (-0.7095,-0.2234) -- (-0.7271,-0.2130) -- (-0.7448,-0.2022) -- (-0.7626,-0.1910) -- (-0.7805,-0.1793) -- (-0.7985,-0.1672) -- (-0.8165,-0.1547) -- (-0.8346,-0.1417) -- (-0.8528,-0.1282) -- (-0.8710,-0.1142) -- (-0.8893,-0.0996) -- (-0.9077,-0.0845) -- (-0.9261,-0.0688) -- (-0.9445,-0.0526) -- (-0.9630,-0.0357) -- (-0.9815,-0.0182) -- (-1.0000,-0.0000);
\draw[mirror] (1.0000,0.0000) -- (0.9815,0.0182) -- (0.9630,0.0357) -- (0.9445,0.0526) -- (0.9261,0.0688) -- (0.9077,0.0845) -- (0.8893,0.0996) -- (0.8710,0.1142) -- (0.8528,0.1282) -- (0.8346,0.1417) -- (0.8165,0.1547) -- (0.7985,0.1672) -- (0.7805,0.1793) -- (0.7626,0.1910) -- (0.7448,0.2022) -- (0.7271,0.2130) -- (0.7095,0.2234) -- (0.6919,0.2334) -- (0.6745,0.2430) -- (0.6571,0.2523) -- (0.6398,0.2612) -- (0.6226,0.2698) -- (0.6055,0.2780) -- (0.5884,0.2860) -- (0.5715,0.2936) -- (0.5546,0.3009) -- (0.5378,0.3080) -- (0.5211,0.3147) -- (0.5045,0.3212) -- (0.4879,0.3274) -- (0.4714,0.3333) -- (0.4550,0.3390) -- (0.4386,0.3445) -- (0.4224,0.3497) -- (0.4061,0.3546) -- (0.3900,0.3594) -- (0.3739,0.3639) -- (0.3579,0.3682) -- (0.3419,0.3723) -- (0.3260,0.3761) -- (0.3101,0.3798) -- (0.2943,0.3833) -- (0.2785,0.3865) -- (0.2628,0.3896) -- (0.2471,0.3925) -- (0.2315,0.3951) -- (0.2159,0.3976) -- (0.2003,0.4000) -- (0.1848,0.4021) -- (0.1693,0.4040) -- (0.1538,0.4058) -- (0.1383,0.4074) -- (0.1229,0.4089) -- (0.1075,0.4101) -- (0.0921,0.4112) -- (0.0767,0.4121) -- (0.0614,0.4129) -- (0.0460,0.4135) -- (0.0307,0.4139) -- (0.0153,0.4141) -- (0.0000,0.4142) -- (-0.0153,0.4141) -- (-0.0307,0.4139) -- (-0.0460,0.4135) -- (-0.0614,0.4129) -- (-0.0767,0.4121) -- (-0.0921,0.4112) -- (-0.1075,0.4101) -- (-0.1229,0.4089) -- (-0.1383,0.4074) -- (-0.1538,0.4058) -- (-0.1693,0.4040) -- (-0.1848,0.4021) -- (-0.2003,0.4000) -- (-0.2159,0.3976) -- (-0.2315,0.3951) -- (-0.2471,0.3925) -- (-0.2628,0.3896) -- (-0.2785,0.3865) -- (-0.2943,0.3833) -- (-0.3101,0.3798) -- (-0.3260,0.3761) -- (-0.3419,0.3723) -- (-0.3579,0.3682) -- (-0.3739,0.3639) -- (-0.3900,0.3594) -- (-0.4061,0.3546) -- (-0.4224,0.3497) -- (-0.4386,0.3445) -- (-0.4550,0.3390) -- (-0.4714,0.3333) -- (-0.4879,0.3274) -- (-0.5045,0.3212) -- (-0.5211,0.3147) -- (-0.5378,0.3080) -- (-0.5546,0.3009) -- (-0.5715,0.2936) -- (-0.5884,0.2860) -- (-0.6055,0.2780) -- (-0.6226,0.2698) -- (-0.6398,0.2612) -- (-0.6571,0.2523) -- (-0.6745,0.2430) -- (-0.6919,0.2334) -- (-0.7095,0.2234) -- (-0.7271,0.2130) -- (-0.7448,0.2022) -- (-0.7626,0.1910) -- (-0.7805,0.1793) -- (-0.7985,0.1672) -- (-0.8165,0.1547) -- (-0.8346,0.1417) -- (-0.8528,0.1282) -- (-0.8710,0.1142) -- (-0.8893,0.0996) -- (-0.9077,0.0845) -- (-0.9261,0.0688) -- (-0.9445,0.0526) -- (-0.9630,0.0357) -- (-0.9815,0.0182) -- (-1.0000,0.0000);
\node[ax4] at (1.0000,0.0000) {};
\node[ax4] at (-1.0000,0.0000) {};
\node[ax4] at (0.0000,0.0000) {};
\node[ax4] at (0.0000,1.0000) {};
\node[ax4] at (0.0000,-1.0000) {};
\node[ax2, rotate=0] at (0.4142,0.0000) {};
\node[ax2, rotate=180] at (-0.4142,0.0000) {};
\node[ax2, rotate=45] at (0.7071,0.7071) {};
\node[ax2, rotate=-45] at (0.7071,-0.7071) {};
\node[ax2, rotate=135] at (-0.7071,0.7071) {};
\node[ax2, rotate=-135] at (-0.7071,-0.7071) {};
\node[ax2, rotate=90] at (0.0000,0.4142) {};
\node[ax2, rotate=-90] at (0.0000,-0.4142) {};
\node[ax3] at (0.3660,0.3660) {};
\node[ax3] at (0.3660,-0.3660) {};
\node[ax3] at (-0.3660,0.3660) {};
\node[ax3] at (-0.3660,-0.3660) {};

\end{tikzpicture}
\caption{Элементы симметрии группы $m\br3m$ ($O_h$) на стереографической
проекции $(001)$ (расчёт): оси 4 (квадраты), 3 (треугольники), 2 (овалы);
линии~— девять плоскостей симметрии (основной круг~— плоскость $(001)$).
Без плоскостей это группа $432$; оставив лишь оси 2 вдоль $\nfam{100}$ и
оси 3, получим группу $23$}
\label{fig:m3m}
\end{figure}

\paragraph{Важные подразделения классов}
\begin{itemize}
\item \emph{Центр симметрии.} В классах с центром симметрии нет полярных
направлений, а значит, нет и полярных свойств~— пироэлектричества,
пьезоэлектричества. Дифракционная картина всегда центросимметрична
(закон Фриделя, ниже), поэтому дифракционными методами различают лишь
11 \term{лауэвских классов}~— центросимметричных:
$\br1$, $2/m$, $mmm$, $4/m$, $4/mmm$, $\br3$, $\br3m$, $6/m$, $6/mmm$, $m\br3$, $m\br3m$.
\item \emph{Энантиоморфизм.} Объекты, симметрия которых содержит только
повороты, могут существовать в двух зеркально равных («правой» и
«левой») формах~— \term{энантиоморфных}; это свойство называют также
\emph{хиральностью} (правый и левый кварц). Энантиоморфных классов 11:
$1$, $2$, $3$, $4$, $6$, $222$, $32$, $422$, $622$, $23$, $432$.
\item \emph{Простые формы.} Многогранник, все грани которого получаются из
одной симметрией класса,~— \term{простая форма}. В классе $m\br3m$ из грани
$(100)$ получается куб, из $(111)$~— октаэдр, из $(110)$~— ромбический
додекаэдр (рис.~\ref{fig:forms}).
\end{itemize}

\begin{figure}[ht]
\centering
\tikzset{bzfront/.style={atomLine, line width=0.7pt, line join=round},
  bzback/.style={atomLine!55, line width=0.4pt, dashed}}
\begin{subfigure}[b]{0.3\textwidth}\centering
\begin{tikzpicture}[cubeproj, scale=1.7]
  \fill[fmBlue!12] (0,0,0) -- (1,0,0) -- (1,0,1) -- (0,0,1) -- cycle;
  \fill[fmBlue!20] (1,0,0) -- (1,1,0) -- (1,1,1) -- (1,0,1) -- cycle;
  \fill[fmBlue!8] (0,0,1) -- (1,0,1) -- (1,1,1) -- (0,1,1) -- cycle;
  \CubeEdges
\end{tikzpicture}
\caption{куб $\{100\}$}
\end{subfigure}\hfill
\begin{subfigure}[b]{0.3\textwidth}\centering
\begin{tikzpicture}[cubeproj, scale=1.0]
  \coordinate (X) at (1,0,0); \coordinate (x) at (-1,0,0);
  \coordinate (Y) at (0,1,0); \coordinate (y) at (0,-1,0);
  \coordinate (Z) at (0,0,1.25); \coordinate (z) at (0,0,-1.25);
  \fill[fmBlue!20] (X) -- (y) -- (Z) -- cycle; \fill[fmBlue!12] (x) -- (y) -- (Z) -- cycle;
  \fill[fmBlue!10] (X) -- (y) -- (z) -- cycle; \fill[fmBlue!6] (x) -- (y) -- (z) -- cycle;
  \draw[atomLine!55, dashed, line width=0.4pt] (Y) -- (X) (Y) -- (x) (Y) -- (Z) (Y) -- (z);
  \draw[atomLine, line width=0.7pt, line join=round] (X) -- (Z) -- (x) -- (z) -- cycle (y) -- (X) (y) -- (Z) (y) -- (x) (y) -- (z);
\end{tikzpicture}
\caption{октаэдр $\{111\}$}
\end{subfigure}\hfill
\begin{subfigure}[b]{0.3\textwidth}\centering
\begin{tikzpicture}[scale=1.25]
  % сгенерировано tools/gen_bands.py — не править вручную
\fill[fmBlue!18] (-0.676,-0.539) -- (-0.469,-0.273) -- (-0.676,0.412) -- (-0.883,0.145) -- cycle;
\fill[fmBlue!14] (0.207,0.267) -- (-0.469,-0.273) -- (0.207,-0.684) -- (0.883,-0.145) -- cycle;
\fill[fmBlue!12] (0.207,-0.684) -- (-0.469,-0.273) -- (-0.676,-0.539) -- (0.000,-0.951) -- cycle;
\fill[fmBlue!20] (-0.676,0.412) -- (-0.469,-0.273) -- (0.207,0.267) -- (0.000,0.951) -- cycle;
\fill[fmBlue!8] (0.676,-0.412) -- (0.883,-0.145) -- (0.207,-0.684) -- (0.000,-0.951) -- cycle;
\fill[fmBlue!16] (0.207,0.267) -- (0.883,-0.145) -- (0.676,0.539) -- (0.000,0.951) -- cycle;
\draw[bzback] (0.469,0.273) -- (-0.207,0.684);
\draw[bzback] (0.469,0.273) -- (-0.207,-0.267);
\draw[bzback] (-0.883,0.145) -- (-0.207,-0.267);
\draw[bzback] (-0.883,0.145) -- (-0.207,0.684);
\draw[bzback] (0.469,0.273) -- (0.676,-0.412);
\draw[bzback] (0.469,0.273) -- (0.676,0.539);
\draw[bzfront] (0.883,-0.145) -- (0.676,0.539);
\draw[bzfront] (0.883,-0.145) -- (0.676,-0.412);
\draw[bzfront] (-0.469,-0.273) -- (-0.676,-0.539);
\draw[bzfront] (-0.469,-0.273) -- (-0.676,0.412);
\draw[bzfront] (-0.883,0.145) -- (-0.676,0.412);
\draw[bzfront] (-0.883,0.145) -- (-0.676,-0.539);
\draw[bzfront] (-0.469,-0.273) -- (0.207,0.267);
\draw[bzfront] (-0.469,-0.273) -- (0.207,-0.684);
\draw[bzfront] (0.883,-0.145) -- (0.207,-0.684);
\draw[bzfront] (0.883,-0.145) -- (0.207,0.267);
\draw[bzfront] (0.000,-0.951) -- (0.676,-0.412);
\draw[bzback] (0.000,-0.951) -- (-0.207,-0.267);
\draw[bzback] (0.000,0.951) -- (-0.207,0.684);
\draw[bzfront] (0.000,0.951) -- (0.676,0.539);
\draw[bzfront] (0.000,-0.951) -- (-0.676,-0.539);
\draw[bzfront] (0.000,-0.951) -- (0.207,-0.684);
\draw[bzfront] (0.000,0.951) -- (0.207,0.267);
\draw[bzfront] (0.000,0.951) -- (-0.676,0.412);
\coordinate (bzG) at (0.000,0.000);
\coordinate (bzH) at (-0.469,-0.273);
\coordinate (bzN) at (0.207,-0.209);
\coordinate (bzP) at (0.207,0.267);

\end{tikzpicture}
\caption{ромбододекаэдр $\{110\}$}
\end{subfigure}
\caption{Простые формы класса $m\br3m$: куб (6 граней), октаэдр (8 граней),
ромбический додекаэдр (12 граней); число граней~— множитель
повторяемости формы}
\label{fig:forms}
\end{figure}

\paragraph{Решётки Браве}
О.~Браве показал, что все многообразие трансляционной симметрии
кристаллов исчерпывается \vazhno{14 типами решёток}, распределёнными по
7 сингониям (рис.~\ref{fig:bravais}, табл.~\ref{tab:bravais}). Ячейку Браве
выбирают по трём условиям, применяемым \emph{последовательно} (первое
важнее второго, второе~— третьего):
\begin{enumerate}
\item симметрия ячейки должна соответствовать симметрии кристалла
(высшей симметрии его сингонии), а рёбра~— быть трансляциями решётки;
\item ячейка должна иметь максимальное число прямых углов, равных углов
и равных рёбер;
\item при этом её объём должен быть минимальным.
\end{enumerate}
По расположению узлов ячейки бывают примитивные ($P$), базоцентрированные
($C$, $A$ или $B$~— узлы в центрах пары параллельных граней),
объёмноцентрированные ($I$~— узел в центре) и гранецентрированные ($F$~—
узлы в центрах всех граней); в тригональной сингонии примитивную
ромбоэдрическую ячейку обозначают $R$.

\begin{table}[ht]
\centering
\caption{Ячейки Браве по сингониям}\label{tab:bravais}
\small
\begin{tabular}{@{}lc@{\qquad}lc@{}}
\toprule
Сингония & Ячейки & Сингония & Ячейки\\
\midrule
Триклинная & $P$ & Тетрагональная & $P$, $I$\\
Моноклинная & $P$, $C$ & Тригональная & $R$\\
Ромбическая & $P$, $C$, $I$, $F$ & Гексагональная & $P$\\
 & & Кубическая & $P$, $I$, $F$\\
\bottomrule
\end{tabular}
\end{table}

\begin{figure}[p]
\centering
\tikzset{bravE/.style={atomLine, line width=0.6pt},
  bravH/.style={atomLine!50, line width=0.45pt, dashed},
  bravN/.style={circle, fill=atomA, draw=atomLine, line width=0.3pt, inner sep=0pt, minimum size=4.2pt},
  bravC/.style={circle, fill=atomB, draw=atomLine, line width=0.3pt, inner sep=0pt, minimum size=4.2pt}}
\begin{tikzpicture}[scale=1.15]
  % создано tools/gen_cryst.py: 14 решёток Браве
\draw[bravE] (0.000,0.000) -- (1.100,0.000);
\draw[bravH] (0.000,0.000) -- (0.720,0.300);
\draw[bravE] (0.000,0.000) -- (0.334,1.160);
\draw[bravE] (0.334,1.160) -- (1.434,1.160);
\draw[bravE] (0.334,1.160) -- (1.054,1.460);
\draw[bravH] (0.720,0.300) -- (1.820,0.300);
\draw[bravH] (0.720,0.300) -- (1.054,1.460);
\draw[bravE] (1.054,1.460) -- (2.154,1.460);
\draw[bravE] (1.100,0.000) -- (1.820,0.300);
\draw[bravE] (1.100,0.000) -- (1.434,1.160);
\draw[bravE] (1.434,1.160) -- (2.154,1.460);
\draw[bravE] (1.820,0.300) -- (2.154,1.460);
\node[bravN] at (0.000,0.000) {};
\node[bravN] at (0.334,1.160) {};
\node[bravN] at (0.720,0.300) {};
\node[bravN] at (1.054,1.460) {};
\node[bravN] at (1.100,0.000) {};
\node[bravN] at (1.434,1.160) {};
\node[bravN] at (1.820,0.300) {};
\node[bravN] at (2.154,1.460) {};
\node[font=\small, anchor=north] at (0.750,-0.250) {триклинная $P$};
\draw[bravE] (3.400,0.000) -- (4.500,0.000);
\draw[bravH] (3.400,0.000) -- (3.820,0.300);
\draw[bravE] (3.400,0.000) -- (3.750,1.050);
\draw[bravE] (3.750,1.050) -- (4.850,1.050);
\draw[bravE] (3.750,1.050) -- (4.170,1.350);
\draw[bravH] (3.820,0.300) -- (4.920,0.300);
\draw[bravH] (3.820,0.300) -- (4.170,1.350);
\draw[bravE] (4.170,1.350) -- (5.270,1.350);
\draw[bravE] (4.500,0.000) -- (4.920,0.300);
\draw[bravE] (4.500,0.000) -- (4.850,1.050);
\draw[bravE] (4.850,1.050) -- (5.270,1.350);
\draw[bravE] (4.920,0.300) -- (5.270,1.350);
\node[bravN] at (3.400,0.000) {};
\node[bravN] at (3.750,1.050) {};
\node[bravN] at (3.820,0.300) {};
\node[bravN] at (4.170,1.350) {};
\node[bravN] at (4.500,0.000) {};
\node[bravN] at (4.850,1.050) {};
\node[bravN] at (4.920,0.300) {};
\node[bravN] at (5.270,1.350) {};
\node[font=\small, anchor=north] at (4.150,-0.250) {моноклинная $P$};
\draw[bravE] (6.800,0.000) -- (7.900,0.000);
\draw[bravH] (6.800,0.000) -- (7.220,0.300);
\draw[bravE] (6.800,0.000) -- (7.150,1.050);
\draw[bravE] (7.150,1.050) -- (8.250,1.050);
\draw[bravE] (7.150,1.050) -- (7.570,1.350);
\draw[bravH] (7.220,0.300) -- (8.320,0.300);
\draw[bravH] (7.220,0.300) -- (7.570,1.350);
\draw[bravE] (7.570,1.350) -- (8.670,1.350);
\draw[bravE] (7.900,0.000) -- (8.320,0.300);
\draw[bravE] (7.900,0.000) -- (8.250,1.050);
\draw[bravE] (8.250,1.050) -- (8.670,1.350);
\draw[bravE] (8.320,0.300) -- (8.670,1.350);
\node[bravN] at (6.800,0.000) {};
\node[bravN] at (7.150,1.050) {};
\node[bravN] at (7.220,0.300) {};
\node[bravN] at (7.570,1.350) {};
\node[bravN] at (7.900,0.000) {};
\node[bravN] at (8.250,1.050) {};
\node[bravN] at (8.320,0.300) {};
\node[bravN] at (8.670,1.350) {};
\node[bravC] at (7.560,0.150) {};
\node[bravC] at (7.910,1.200) {};
\node[font=\small, anchor=north] at (7.550,-0.250) {моноклинная $C$};
\draw[bravE] (10.200,0.000) -- (10.304,0.664);
\draw[bravE] (10.200,0.000) -- (10.858,1.059);
\draw[bravE] (10.200,0.000) -- (9.439,0.861);
\draw[bravE] (9.439,0.861) -- (9.542,1.525);
\draw[bravH] (9.439,0.861) -- (10.096,1.920);
\draw[bravE] (10.858,1.059) -- (10.961,1.723);
\draw[bravH] (10.858,1.059) -- (10.096,1.920);
\draw[bravH] (10.096,1.920) -- (10.200,2.584);
\draw[bravE] (10.304,0.664) -- (10.961,1.723);
\draw[bravE] (10.304,0.664) -- (9.542,1.525);
\draw[bravE] (9.542,1.525) -- (10.200,2.584);
\draw[bravE] (10.961,1.723) -- (10.200,2.584);
\node[bravN] at (10.200,0.000) {};
\node[bravN] at (9.439,0.861) {};
\node[bravN] at (10.858,1.059) {};
\node[bravN] at (10.096,1.920) {};
\node[bravN] at (10.304,0.664) {};
\node[bravN] at (9.542,1.525) {};
\node[bravN] at (10.961,1.723) {};
\node[bravN] at (10.200,2.584) {};
\node[font=\small, anchor=north] at (10.950,-0.250) {тригональная $R$};
\draw[bravE] (0.000,-3.050) -- (1.250,-3.050);
\draw[bravH] (0.000,-3.050) -- (0.378,-2.780);
\draw[bravE] (0.000,-3.050) -- (0.000,-2.050);
\draw[bravE] (0.000,-2.050) -- (1.250,-2.050);
\draw[bravE] (0.000,-2.050) -- (0.378,-1.780);
\draw[bravH] (0.378,-2.780) -- (1.628,-2.780);
\draw[bravH] (0.378,-2.780) -- (0.378,-1.780);
\draw[bravE] (0.378,-1.780) -- (1.628,-1.780);
\draw[bravE] (1.250,-3.050) -- (1.628,-2.780);
\draw[bravE] (1.250,-3.050) -- (1.250,-2.050);
\draw[bravE] (1.250,-2.050) -- (1.628,-1.780);
\draw[bravE] (1.628,-2.780) -- (1.628,-1.780);
\node[bravN] at (0.000,-3.050) {};
\node[bravN] at (0.000,-2.050) {};
\node[bravN] at (0.378,-2.780) {};
\node[bravN] at (0.378,-1.780) {};
\node[bravN] at (1.250,-3.050) {};
\node[bravN] at (1.250,-2.050) {};
\node[bravN] at (1.628,-2.780) {};
\node[bravN] at (1.628,-1.780) {};
\node[font=\small, anchor=north] at (0.750,-3.300) {ромбическая $P$};
\draw[bravE] (3.400,-3.050) -- (4.650,-3.050);
\draw[bravH] (3.400,-3.050) -- (3.778,-2.780);
\draw[bravE] (3.400,-3.050) -- (3.400,-2.050);
\draw[bravE] (3.400,-2.050) -- (4.650,-2.050);
\draw[bravE] (3.400,-2.050) -- (3.778,-1.780);
\draw[bravH] (3.778,-2.780) -- (5.028,-2.780);
\draw[bravH] (3.778,-2.780) -- (3.778,-1.780);
\draw[bravE] (3.778,-1.780) -- (5.028,-1.780);
\draw[bravE] (4.650,-3.050) -- (5.028,-2.780);
\draw[bravE] (4.650,-3.050) -- (4.650,-2.050);
\draw[bravE] (4.650,-2.050) -- (5.028,-1.780);
\draw[bravE] (5.028,-2.780) -- (5.028,-1.780);
\node[bravN] at (3.400,-3.050) {};
\node[bravN] at (3.400,-2.050) {};
\node[bravN] at (3.778,-2.780) {};
\node[bravN] at (3.778,-1.780) {};
\node[bravN] at (4.650,-3.050) {};
\node[bravN] at (4.650,-2.050) {};
\node[bravN] at (5.028,-2.780) {};
\node[bravN] at (5.028,-1.780) {};
\node[bravC] at (4.214,-2.915) {};
\node[bravC] at (4.214,-1.915) {};
\node[font=\small, anchor=north] at (4.150,-3.300) {ромбическая $C$};
\draw[bravE] (6.800,-3.050) -- (8.050,-3.050);
\draw[bravH] (6.800,-3.050) -- (7.178,-2.780);
\draw[bravE] (6.800,-3.050) -- (6.800,-2.050);
\draw[bravE] (6.800,-2.050) -- (8.050,-2.050);
\draw[bravE] (6.800,-2.050) -- (7.178,-1.780);
\draw[bravH] (7.178,-2.780) -- (8.428,-2.780);
\draw[bravH] (7.178,-2.780) -- (7.178,-1.780);
\draw[bravE] (7.178,-1.780) -- (8.428,-1.780);
\draw[bravE] (8.050,-3.050) -- (8.428,-2.780);
\draw[bravE] (8.050,-3.050) -- (8.050,-2.050);
\draw[bravE] (8.050,-2.050) -- (8.428,-1.780);
\draw[bravE] (8.428,-2.780) -- (8.428,-1.780);
\node[bravN] at (6.800,-3.050) {};
\node[bravN] at (6.800,-2.050) {};
\node[bravN] at (7.178,-2.780) {};
\node[bravN] at (7.178,-1.780) {};
\node[bravN] at (8.050,-3.050) {};
\node[bravN] at (8.050,-2.050) {};
\node[bravN] at (8.428,-2.780) {};
\node[bravN] at (8.428,-1.780) {};
\node[bravC] at (7.614,-2.415) {};
\node[font=\small, anchor=north] at (7.550,-3.300) {ромбическая $I$};
\draw[bravE] (10.200,-3.050) -- (11.450,-3.050);
\draw[bravH] (10.200,-3.050) -- (10.578,-2.780);
\draw[bravE] (10.200,-3.050) -- (10.200,-2.050);
\draw[bravE] (10.200,-2.050) -- (11.450,-2.050);
\draw[bravE] (10.200,-2.050) -- (10.578,-1.780);
\draw[bravH] (10.578,-2.780) -- (11.828,-2.780);
\draw[bravH] (10.578,-2.780) -- (10.578,-1.780);
\draw[bravE] (10.578,-1.780) -- (11.828,-1.780);
\draw[bravE] (11.450,-3.050) -- (11.828,-2.780);
\draw[bravE] (11.450,-3.050) -- (11.450,-2.050);
\draw[bravE] (11.450,-2.050) -- (11.828,-1.780);
\draw[bravE] (11.828,-2.780) -- (11.828,-1.780);
\node[bravN] at (10.200,-3.050) {};
\node[bravN] at (10.200,-2.050) {};
\node[bravN] at (10.578,-2.780) {};
\node[bravN] at (10.578,-1.780) {};
\node[bravN] at (11.450,-3.050) {};
\node[bravN] at (11.450,-2.050) {};
\node[bravN] at (11.828,-2.780) {};
\node[bravN] at (11.828,-1.780) {};
\node[bravC] at (11.014,-2.915) {};
\node[bravC] at (11.014,-1.915) {};
\node[bravC] at (10.389,-2.415) {};
\node[bravC] at (11.639,-2.415) {};
\node[bravC] at (10.825,-2.550) {};
\node[bravC] at (11.203,-2.280) {};
\node[font=\small, anchor=north] at (10.950,-3.300) {ромбическая $F$};
\draw[bravE] (0.000,-6.100) -- (0.950,-6.100);
\draw[bravH] (0.000,-6.100) -- (0.399,-5.815);
\draw[bravE] (0.000,-6.100) -- (0.000,-4.750);
\draw[bravE] (0.000,-4.750) -- (0.950,-4.750);
\draw[bravE] (0.000,-4.750) -- (0.399,-4.465);
\draw[bravH] (0.399,-5.815) -- (1.349,-5.815);
\draw[bravH] (0.399,-5.815) -- (0.399,-4.465);
\draw[bravE] (0.399,-4.465) -- (1.349,-4.465);
\draw[bravE] (0.950,-6.100) -- (1.349,-5.815);
\draw[bravE] (0.950,-6.100) -- (0.950,-4.750);
\draw[bravE] (0.950,-4.750) -- (1.349,-4.465);
\draw[bravE] (1.349,-5.815) -- (1.349,-4.465);
\node[bravN] at (0.000,-6.100) {};
\node[bravN] at (0.000,-4.750) {};
\node[bravN] at (0.399,-5.815) {};
\node[bravN] at (0.399,-4.465) {};
\node[bravN] at (0.950,-6.100) {};
\node[bravN] at (0.950,-4.750) {};
\node[bravN] at (1.349,-5.815) {};
\node[bravN] at (1.349,-4.465) {};
\node[font=\small, anchor=north] at (0.750,-6.350) {тетрагональная $P$};
\draw[bravE] (3.400,-6.100) -- (4.350,-6.100);
\draw[bravH] (3.400,-6.100) -- (3.799,-5.815);
\draw[bravE] (3.400,-6.100) -- (3.400,-4.750);
\draw[bravE] (3.400,-4.750) -- (4.350,-4.750);
\draw[bravE] (3.400,-4.750) -- (3.799,-4.465);
\draw[bravH] (3.799,-5.815) -- (4.749,-5.815);
\draw[bravH] (3.799,-5.815) -- (3.799,-4.465);
\draw[bravE] (3.799,-4.465) -- (4.749,-4.465);
\draw[bravE] (4.350,-6.100) -- (4.749,-5.815);
\draw[bravE] (4.350,-6.100) -- (4.350,-4.750);
\draw[bravE] (4.350,-4.750) -- (4.749,-4.465);
\draw[bravE] (4.749,-5.815) -- (4.749,-4.465);
\node[bravN] at (3.400,-6.100) {};
\node[bravN] at (3.400,-4.750) {};
\node[bravN] at (3.799,-5.815) {};
\node[bravN] at (3.799,-4.465) {};
\node[bravN] at (4.350,-6.100) {};
\node[bravN] at (4.350,-4.750) {};
\node[bravN] at (4.749,-5.815) {};
\node[bravN] at (4.749,-4.465) {};
\node[bravC] at (4.074,-5.282) {};
\node[font=\small, anchor=north] at (4.150,-6.350) {тетрагональная $I$};
\draw[bravE] (6.800,-6.100) -- (7.900,-6.100);
\draw[bravH] (6.800,-6.100) -- (6.650,-5.814);
\draw[bravE] (6.800,-6.100) -- (6.800,-4.850);
\draw[bravE] (6.800,-4.850) -- (7.900,-4.850);
\draw[bravE] (6.800,-4.850) -- (6.650,-4.564);
\draw[bravH] (6.650,-5.814) -- (7.750,-5.814);
\draw[bravH] (6.650,-5.814) -- (6.650,-4.564);
\draw[bravE] (6.650,-4.564) -- (7.750,-4.564);
\draw[bravE] (7.900,-6.100) -- (7.750,-5.814);
\draw[bravE] (7.900,-6.100) -- (7.900,-4.850);
\draw[bravE] (7.900,-4.850) -- (7.750,-4.564);
\draw[bravE] (7.750,-5.814) -- (7.750,-4.564);
\node[bravN] at (6.800,-6.100) {};
\node[bravN] at (6.800,-4.850) {};
\node[bravN] at (6.650,-5.814) {};
\node[bravN] at (6.650,-4.564) {};
\node[bravN] at (7.900,-6.100) {};
\node[bravN] at (7.900,-4.850) {};
\node[bravN] at (7.750,-5.814) {};
\node[bravN] at (7.750,-4.564) {};
\node[font=\small, anchor=north] at (7.550,-6.350) {гексагональная $P$};
\draw[bravE] (0.000,-9.150) -- (1.100,-9.150);
\draw[bravH] (0.000,-9.150) -- (0.462,-8.820);
\draw[bravE] (0.000,-9.150) -- (0.000,-8.050);
\draw[bravE] (0.000,-8.050) -- (1.100,-8.050);
\draw[bravE] (0.000,-8.050) -- (0.462,-7.720);
\draw[bravH] (0.462,-8.820) -- (1.562,-8.820);
\draw[bravH] (0.462,-8.820) -- (0.462,-7.720);
\draw[bravE] (0.462,-7.720) -- (1.562,-7.720);
\draw[bravE] (1.100,-9.150) -- (1.562,-8.820);
\draw[bravE] (1.100,-9.150) -- (1.100,-8.050);
\draw[bravE] (1.100,-8.050) -- (1.562,-7.720);
\draw[bravE] (1.562,-8.820) -- (1.562,-7.720);
\node[bravN] at (0.000,-9.150) {};
\node[bravN] at (0.000,-8.050) {};
\node[bravN] at (0.462,-8.820) {};
\node[bravN] at (0.462,-7.720) {};
\node[bravN] at (1.100,-9.150) {};
\node[bravN] at (1.100,-8.050) {};
\node[bravN] at (1.562,-8.820) {};
\node[bravN] at (1.562,-7.720) {};
\node[font=\small, anchor=north] at (0.750,-9.400) {кубическая $P$};
\draw[bravE] (3.400,-9.150) -- (4.500,-9.150);
\draw[bravH] (3.400,-9.150) -- (3.862,-8.820);
\draw[bravE] (3.400,-9.150) -- (3.400,-8.050);
\draw[bravE] (3.400,-8.050) -- (4.500,-8.050);
\draw[bravE] (3.400,-8.050) -- (3.862,-7.720);
\draw[bravH] (3.862,-8.820) -- (4.962,-8.820);
\draw[bravH] (3.862,-8.820) -- (3.862,-7.720);
\draw[bravE] (3.862,-7.720) -- (4.962,-7.720);
\draw[bravE] (4.500,-9.150) -- (4.962,-8.820);
\draw[bravE] (4.500,-9.150) -- (4.500,-8.050);
\draw[bravE] (4.500,-8.050) -- (4.962,-7.720);
\draw[bravE] (4.962,-8.820) -- (4.962,-7.720);
\node[bravN] at (3.400,-9.150) {};
\node[bravN] at (3.400,-8.050) {};
\node[bravN] at (3.862,-8.820) {};
\node[bravN] at (3.862,-7.720) {};
\node[bravN] at (4.500,-9.150) {};
\node[bravN] at (4.500,-8.050) {};
\node[bravN] at (4.962,-8.820) {};
\node[bravN] at (4.962,-7.720) {};
\node[bravC] at (4.181,-8.435) {};
\node[font=\small, anchor=north] at (4.150,-9.400) {кубическая $I$};
\draw[bravE] (6.800,-9.150) -- (7.900,-9.150);
\draw[bravH] (6.800,-9.150) -- (7.262,-8.820);
\draw[bravE] (6.800,-9.150) -- (6.800,-8.050);
\draw[bravE] (6.800,-8.050) -- (7.900,-8.050);
\draw[bravE] (6.800,-8.050) -- (7.262,-7.720);
\draw[bravH] (7.262,-8.820) -- (8.362,-8.820);
\draw[bravH] (7.262,-8.820) -- (7.262,-7.720);
\draw[bravE] (7.262,-7.720) -- (8.362,-7.720);
\draw[bravE] (7.900,-9.150) -- (8.362,-8.820);
\draw[bravE] (7.900,-9.150) -- (7.900,-8.050);
\draw[bravE] (7.900,-8.050) -- (8.362,-7.720);
\draw[bravE] (8.362,-8.820) -- (8.362,-7.720);
\node[bravN] at (6.800,-9.150) {};
\node[bravN] at (6.800,-8.050) {};
\node[bravN] at (7.262,-8.820) {};
\node[bravN] at (7.262,-7.720) {};
\node[bravN] at (7.900,-9.150) {};
\node[bravN] at (7.900,-8.050) {};
\node[bravN] at (8.362,-8.820) {};
\node[bravN] at (8.362,-7.720) {};
\node[bravC] at (7.581,-8.985) {};
\node[bravC] at (7.581,-7.885) {};
\node[bravC] at (7.031,-8.435) {};
\node[bravC] at (8.131,-8.435) {};
\node[bravC] at (7.350,-8.600) {};
\node[bravC] at (7.812,-8.270) {};
\node[font=\small, anchor=north] at (7.550,-9.400) {кубическая $F$};

\end{tikzpicture}
\caption{Четырнадцать решёток Браве (расчёт). Оранжевым выделены
центрирующие узлы. Почему, например, нет тетрагональной $C$-решётки?
Потому что она сводится к примитивной тетрагональной $P$ с меньшей
ячейкой, повёрнутой на $45^\circ$; по той же причине тетрагональная $F$
сводится к $I$}
\label{fig:bravais}
\end{figure}

\subsection{Пространственные группы}\label{ssec:space-groups}

\begin{opredelenie}[Пространственная группа]
Совокупность всех преобразований кристаллического пространства,
переводящих каждую точку и каждое направление в эквивалентные.
Элемент группы~— преобразование
\begin{equation}\label{eq:seitz}
\vb r'=R\vb r+\vb t\equiv\{R\,|\,\vb t\}\,\vb r ,
\end{equation}
где $R$~— поворот (собственный или несобственный), $\vb t$~— перенос.
\end{opredelenie}

При $\det R=+1$ и $\vb t\ne0$ получается винтовое движение, при $\det R=-1$~—
отражение со сдвигом (скользящее отражение). Последовательное
применение двух операций:
\begin{equation}\label{eq:seitz-product}
\{R_2|\vb t_2\}\{R_1|\vb t_1\}\vb r=R_2(R_1\vb r+\vb t_1)+\vb t_2
\ \Rightarrow\ \{R_2|\vb t_2\}\{R_1|\vb t_1\}=\{R_2R_1\,|\,R_2\vb t_1+\vb t_2\}
\end{equation}
\ispr{В оригинале результат записан как «$R_2R_1+R_2\vb t_1+\vb t_2$», что
смешивает матрицу и вектор; произведение~— снова операция вида
$\{R|\vb t\}$ с $R=R_2R_1$ и $\vb t=R_2\vb t_1+\vb t_2$.}; единичный элемент~—
$\{E|\vb 0\}$, обратный~— $\{R|\vb t\}^{-1}=\{R^{-1}|-R^{-1}\vb t\}$. Операции $\{E|\vb t\}$
с векторами решётки образуют бесконечную \emph{трансляционную} подгруппу;
всего таких групп 14~— по числу решёток Браве.

\paragraph{Винтовые оси}
\term{Винтовая ось}~— ось симметрии и параллельный ей перенос,
действующие совместно. Обозначение $n_m$: поворот на $360^\circ/n$ и сдвиг на
$m/n$ периода вдоль оси. Например, $3_1$~— поворот на $120^\circ$ и сдвиг на
$\frac13$ периода, $3_2$~— сдвиг на $\frac23$ (рис.~\ref{fig:screw}). Оси $3_1$ и $3_2$
(так же как $4_1$ и $4_3$, $6_1$ и $6_5$)~— правая и левая разновидности, они
энантиоморфны; ось $2_1$ можно считать и правой, и левой. Действие оси
$2_1$ вдоль $\vb a_2$:
\[
\begin{pmatrix}x'\\y'\\z'\end{pmatrix}=
\begin{pmatrix}\br1&0&0\\0&1&0\\0&0&\br1\end{pmatrix}\begin{pmatrix}x\\y\\z\end{pmatrix}+
\begin{pmatrix}0\\\frac12\\0\end{pmatrix}=\begin{pmatrix}-x\\y+\frac12\\-z\end{pmatrix};
\]
двукратное применение даёт $(x,\ y+1,\ z)$~— трансляцию на период.

\paragraph{Плоскости скользящего отражения}
\term{Плоскость скользящего отражения}~— отражение в плоскости и
параллельный ей перенос, действующие совместно (рис.~\ref{fig:glide}).
Различают плоскости $a$, $b$, $c$ (сдвиг на половину периода вдоль
соответствующей оси), $n$ (сдвиг вдоль диагонали грани:
$(\vb a+\vb b)/2$ и т.\,п.) и «алмазную» $d$ (сдвиг на $(\vb a\pm\vb b)/4$,
$(\vb a\pm\vb b\pm\vb c)/4$ и т.\,п.). Для плоскости, перпендикулярной $\vb a_2$, со
сдвигом вдоль $\vb a_3$: $(x,y,z)\to(x,-y,z+\frac12)$, а двукратное применение
снова даёт трансляцию.

\begin{figure}[ht]
\centering
\begin{subfigure}[b]{0.56\textwidth}\centering
\begin{tikzpicture}[scale=1.0]
  \foreach \col/\lab/\m in {0/3/0, 2.3/3_1/1, 4.6/3_2/2}{
    \draw[fmGray, thick] (\col,-0.2) -- (\col,3.3);
    \foreach \k in {0,1,2,3}{
      \pgfmathsetmacro{\ang}{120*\k}
      \pgfmathsetmacro{\hz}{\m==0 ? 0 : mod(\m*\k,3)}
      \pgfmathsetmacro{\hh}{\m==0 ? 0 : \hz + (\k==3 ? 3 : 0)}
      \pgfmathsetmacro{\px}{\col+0.65*cos(\ang)}
      \pgfmathsetmacro{\py}{\hh+0.22*sin(\ang)}
      \ifnum\k<3
        \fill[fmBlue] (\px,\py) circle (0.08);
        \node[font=\tiny, anchor=west] at (\px+0.07,\py) {\pgfmathparse{int(\k+1)}\pgfmathresult};
      \fi}
    \ifnum\m>0
      \foreach \k in {0,1,2}{
        \pgfmathsetmacro{\ang}{120*\k}
        \pgfmathsetmacro{\py}{mod(\m*\k,3)+3+0.22*sin(\ang)}
        \pgfmathsetmacro{\px}{\col+0.65*cos(\ang)}
        \ifdim\py pt<3.35pt \fill[fmBlue!40] (\px,\py) circle (0.08);\fi}
    \fi
    \draw[fmGray, densely dotted] (\col-0.75,0) -- (\col+0.75,0);
    \draw[fmGray, densely dotted] (\col-0.75,3) -- (\col+0.75,3);
    \node[font=\small] at (\col,-0.55) {$\lab$};}
  \node[font=\scriptsize, anchor=east] at (-0.8,0) {$0$};
  \node[font=\scriptsize, anchor=east] at (-0.8,1) {$\frac13$};
  \node[font=\scriptsize, anchor=east] at (-0.8,2) {$\frac23$};
  \node[font=\scriptsize, anchor=east] at (-0.8,3) {$1$};
  \foreach \yy in {1,2}{\draw[fmGray!50, densely dotted] (-0.75,\yy) -- (5.35,\yy);}
\end{tikzpicture}
\caption{}
\end{subfigure}\hfill
\begin{subfigure}[b]{0.4\textwidth}\centering
\begin{tikzpicture}[scale=0.95]
  \draw[fmRed, thick, dashed] (-0.3,0) -- (5.0,0) node[right, font=\small]{$c$};
  \foreach \i in {0,1,2}{
    \begin{scope}[xshift=\i*1.8cm]
      \draw[fmBlue, thick, fill=fmBlue!20] (0,0.35) -- (0,1.0) -- (0.45,1.0) -- (0.45,0.85) -- (0.15,0.85) -- (0.15,0.35) -- cycle;
    \end{scope}
    \begin{scope}[xshift=\i*1.8cm+0.9cm]
      \draw[fmGreen!60!black, thick, fill=fmGreen!20] (0,-0.35) -- (0,-1.0) -- (0.45,-1.0) -- (0.45,-0.85) -- (0.15,-0.85) -- (0.15,-0.35) -- cycle;
    \end{scope}}
  \draw[dimline] (0,-1.35) -- node[below, font=\scriptsize]{$\vb t=\vb c/2$} (0.9,-1.35);
\end{tikzpicture}
\caption{}
\end{subfigure}
\caption{Открытые элементы симметрии: \textit{а}~— поворотная ось 3 и
винтовые оси $3_1$ и $3_2$ (вид сбоку; числа~— порядок, в котором точки
получаются из первой; бледные~— точки следующего периода); \textit{б}~—
плоскость скользящего отражения $c$ (штрих): фигура отражается в
плоскости и сдвигается на половину периода}
\label{fig:screw}\label{fig:glide}
\end{figure}

\paragraph{230 пространственных групп}
Перебрав все сочетания открытых и закрытых элементов симметрии,
Е.\,С.~Фёдоров и А.~Шёнфлис (независимо, 1890--1891) получили \vazhno{230
пространственных групп}. В 73 \emph{симморфных} группах точечная группа
является подгруппой пространственной; в остальных (несимморфных) этого
нет из-за винтовых осей и плоскостей скольжения.

\paragraph{Правильные системы точек}
Совокупность точек, получаемых из одной всеми операциями
пространственной группы,~— \term{правильная система точек} (аналог
простой формы для точечной группы). Если исходная точка лежит на
каком-либо элементе симметрии (или на равных расстояниях от одинаковых
элементов), система \emph{частная}, иначе~— \emph{общая}.
\term{Кратность}~— число точек системы в элементарной ячейке; у общей
системы она больше, чем у частных. Атомы одного сорта в структуре
занимают одну или несколько правильных систем точек.

\paragraph{Международный символ пространственной группы}
Первая буква~— тип решётки Браве; далее~— порождающие элементы
симметрии, каждый на своём месте (перестановка меняет смысл). В
кубической сингонии на третьем месте всегда стоит 3 (оси 3 вдоль
$\nfam{111}$), на втором~— элементы, связанные с $\nfam{100}$, на
четвёртом~— с $\nfam{110}$. Примеры: медь~— $Fm\br3m$, вольфрам~— $Im\br3m$,
алмаз~— $Fd\br3m$ (вдоль $\nfam{100}$ проходит алмазная плоскость $d$),
NaCl~— $Fm\br3m$, CsCl~— $Pm\br3m$.

\paragraph{Антисимметрия и цветная симметрия}
А.\,В.~Шубников ввёл операцию \emph{антисимметрии}, меняющую «знак»
объекта (как электрон и позитрон, негатив и позитив). Чёрно-белые группы
хорошо описывают магнитные структуры (направление спина~— «знак»): к
32 обычным точечным группам добавляется 58 чёрно-белых~— всего 90
магнитных классов; пространственных шубниковских групп 1651. Обобщение на
многоцветные группы предложил Н.\,В.~Белов.

\paragraph{Двумерная и многомерная кристаллография. Квазикристаллы}
На плоскости возможны те же оси 1, 2, 3, 4, 6; сингоний четыре
(косоугольная, прямоугольная, квадратная, гексагональная), решёток Браве
пять (в прямоугольной есть центрированная), плоских групп симметрии~—
17. В четырёхмерной кристаллографии разрешены оси 5, 8, 10 и 12-го
порядков; решёток Браве там 64, точечных групп 227, пространственных
около 4\,800.

Упорядоченные структуры делят на \emph{периодические} (кристаллы) и
\emph{квазипериодические}, у которых обратная решётка квазипериодична.
Среди последних~— \emph{несоразмерные} кристаллы (кристаллографическая
точечная группа) и \term{квазикристаллы} (некристаллографическая точечная
группа, например с осями 5-го порядка). Двумерный пример~— мозаика Пенроуза;
трёхмерный~— икосаэдрические сплавы Al--Mn (Д.~Шехтман, 1984;
Нобелевская премия 2011~г.).

\paragraph{Как симметрию определяют на опыте}
Атомное строение изучают по дифракции рентгеновских лучей, электронов и
нейтронов: длины волн сравнимы с межатомными расстояниями, и кристалл
служит дифракционной решёткой. Максимумы возникают в направлениях,
удовлетворяющих условию Вульфа--Брэгга
\begin{keyeq}\label{eq:bragg}
2d\sin\theta=n\lambda ,
\end{keyeq}
где $d$~— межплоскостное расстояние, $\theta$~— угол скольжения, $n$~— порядок
отражения, $\lambda$~— длина волны.
\begin{itemize}
\item \emph{Метод Лауэ} (неподвижный монокристалл, «белое» излучение):
определяют ориентацию и элементы симметрии; по совпадению точек первого и
второго рода~— сингонию, по симметрии лауэграммы~— порядок осей.
\item По \emph{закону Фриделя} интенсивности отражений $(hkl)$ и
$(\br h\br k\br l)$ равны: дифракционная картина имеет центр симметрии, даже
если у кристалла его нет. Поэтому классы, различающиеся только центром
симметрии, дают одинаковую картину~— отсюда 11 лауэвских классов.
\item \emph{Метод вращения} (монохроматический пучок, кристалл вращается
вокруг оси зоны): находят период идентичности вдоль оси; сравнивая периоды
вдоль $\nfam{100}$, $\nfam{110}$, $\nfam{111}$, определяют тип решётки Браве.
\item \emph{Метод Дебая--Шеррера}~— для поликристаллов (монохроматическое
излучение).
\end{itemize}
Сначала по геометрии дифракционной картины определяют решётку, затем по
интенсивностям~— расположение атомов в ячейке. По диффузному рассеянию
изучают также упругие волны, модули упругости, фононный спектр.

\subsection{Предельные группы симметрии}\label{ssec:curie}

Для физических полей, тел вращения и многих свойств кристаллов нужны
группы с осями \emph{бесконечного} порядка (поворот на любой угол).
Это \term{предельные группы}, или группы Кюри; П.~Кюри показал, что их
семь (рис.~\ref{fig:curie}, табл.~\ref{tab:curie}). Все 32 кристаллографические
точечные группы~— подгруппы предельных.

\begin{figure}[ht]
\centering
\tikzset{cbody/.style={fill=fmBlue!12, draw=atomLine, line width=0.6pt},
  crot/.style={-{Stealth[length=1.6mm]}, fmRed, line width=0.7pt}}
\begin{tikzpicture}[scale=0.82]
  % а: вращающийся конус
  \begin{scope}[xshift=0cm]
    \draw[cbody] (-0.7,0) -- (0,2.0) -- (0.7,0) arc (0:-180:0.7 and 0.22);
    \draw[atomLine, line width=0.6pt, dashed] (0.7,0) arc (0:180:0.7 and 0.22);
    \draw[crot] (-0.45,1.25) arc (190:520:0.45 and 0.13);
    \node[font=\small] at (0,-0.75) {\textit{а}~$\infty$};
  \end{scope}
  % б: неподвижный конус
  \begin{scope}[xshift=2.2cm]
    \draw[cbody] (-0.7,0) -- (0,2.0) -- (0.7,0) arc (0:-180:0.7 and 0.22);
    \draw[atomLine, line width=0.6pt, dashed] (0.7,0) arc (0:180:0.7 and 0.22);
    \node[font=\small] at (0,-0.75) {\textit{б}~$\infty m$};
  \end{scope}
  % в: вращающийся цилиндр
  \begin{scope}[xshift=4.4cm]
    \draw[cbody] (-0.6,0) -- (-0.6,1.8) arc (180:360:0.6 and 0.18) -- (0.6,0) arc (0:-180:0.6 and 0.18);
    \draw[cbody] (0,1.8) ellipse (0.6 and 0.18);
    \draw[crot] (-0.6,0.9) arc (180:500:0.6 and 0.18);
    \node[font=\small] at (0,-0.75) {\textit{в}~$\infty/m$};
  \end{scope}
  % г: закрученный цилиндр
  \begin{scope}[xshift=6.6cm]
    \draw[cbody] (-0.6,0) -- (-0.6,1.8) arc (180:360:0.6 and 0.18) -- (0.6,0) arc (0:-180:0.6 and 0.18);
    \draw[cbody] (0,1.8) ellipse (0.6 and 0.18);
    \foreach \x in {-0.4,0,0.4}{\draw[fmRed, line width=0.5pt] (\x-0.15,0) .. controls (\x,0.6) and (\x,1.2) .. (\x+0.15,1.8);}
    \node[font=\small] at (0,-0.75) {\textit{г}~$\infty2$};
  \end{scope}
  % д: неподвижный цилиндр
  \begin{scope}[xshift=8.8cm]
    \draw[cbody] (-0.6,0) -- (-0.6,1.8) arc (180:360:0.6 and 0.18) -- (0.6,0) arc (0:-180:0.6 and 0.18);
    \draw[cbody] (0,1.8) ellipse (0.6 and 0.18);
    \node[font=\small] at (0,-0.75) {\textit{д}~$\infty/mmm$};
  \end{scope}
  % е: закрученный шар
  \begin{scope}[xshift=11.0cm, yshift=0.9cm]
    \shade[ball color=fmBlue!25] (0,0) circle (0.85);
    \draw[atomLine, line width=0.6pt] (0,0) circle (0.85);
    \foreach \a in {-50,0,50}{\draw[fmRed, line width=0.5pt, rotate=\a] (-0.6,-0.6) .. controls (-0.2,0.1) and (0.2,-0.1) .. (0.6,0.6);}
    \node[font=\small] at (0,-1.65) {\textit{е}~$\infty\infty$};
  \end{scope}
  % ж: шар
  \begin{scope}[xshift=13.2cm, yshift=0.9cm]
    \shade[ball color=fmBlue!25] (0,0) circle (0.85);
    \draw[atomLine, line width=0.6pt] (0,0) circle (0.85);
    \node[font=\small] at (0,-1.65) {\textit{ж}~$\infty\infty m$};
  \end{scope}
\end{tikzpicture}
\caption{Фигуры, изображающие предельные группы Кюри: \textit{а}~—
вращающийся конус; \textit{б}~— неподвижный конус; \textit{в}~— вращающийся
цилиндр; \textit{г}~— цилиндр с закрученными в разные стороны концами;
\textit{д}~— неподвижный цилиндр; \textit{е}~— шар с закрученными
диаметрами; \textit{ж}~— обычный шар}
\label{fig:curie}
\end{figure}

\begin{table}[ht]
\centering
\caption{Предельные группы симметрии и их физические реализации}\label{tab:curie}
\small
\begin{tabularx}{\textwidth}{@{}lp{30mm}>{\raggedright\arraybackslash}p{28mm}>{\raggedright\arraybackslash\hspace{0pt}}X@{}}
\toprule
Группа & Фигура & Свойства & Пример \\
\midrule
$\infty$ & вращающийся конус & полярна, энантиоморфна & предел групп 1, 2, 3, 4, 6 \\
$\infty m$ & неподвижный конус & полярна & однородное электрическое поле (полярный вектор $\vb E$) \\
$\infty/m$ & вращающийся цилиндр & центр симметрии & магнитное поле (аксиальный вектор $\vb H$) \\
$\infty2$ & закрученный цилиндр & энантиоморфна & вращение плоскости поляризации в анизотропной среде \\
$\infty/mmm$ & неподвижный цилиндр & центр симметрии & одноосное растяжение или сжатие \\
$\infty\infty$ & закрученный шар & энантиоморфна & вращение плоскости поляризации в изотропной среде \\
$\infty\infty m$ & шар & все операции & гидростатическое давление, однородный нагрев (скаляр) \\
\bottomrule
\end{tabularx}
\end{table}

Группа $\infty\infty m$ содержит все точечные операции, оставляющие на месте
шар,~— её называют \emph{ортогональной} группой. Предельные группы особенно
полезны в кристаллофизике: они описывают симметрию внешних воздействий
(п.~\ref{ssec:symmetry-principle}).

\begin{itogi}
\begin{itemize}
\item Трансляционная симметрия допускает только оси 1, 2, 3, 4, 6
($\operatorname{tr}R=1+2\cos\alpha$~— целое).
\item Инверсионные оси: $\br1$~— центр, $\br2=m$~— плоскость, $\br3=3+\br1$,
$\br6=3/m$, $\br4$~— независимый элемент.
\item Группа: замкнутость, ассоциативность, единица, обратный элемент;
подгруппы, классы, генераторы, таблица Кэли.
\item 32 точечные группы (11 лауэвских, 11 энантиоморфных); 14 решёток
Браве; 230 пространственных групп (винтовые оси $n_m$, плоскости
скольжения $a$, $b$, $c$, $n$, $d$).
\item 7 предельных групп Кюри описывают симметрию полей и воздействий:
электрическое поле~— $\infty m$, магнитное~— $\infty/m$, одноосное напряжение~—
$\infty/mmm$.
\end{itemize}
\end{itogi}

\begin{proverka}
\item Почему в кристаллах нет оси 5-го порядка, а в квазикристаллах она
есть?
\item Чему равен след матрицы поворота вокруг оси 6? Оси $\br4$?
\item Покажите, что ось $\br2$ эквивалентна плоскости симметрии, а $\br6$~—
сочетанию оси 3 и перпендикулярной плоскости.
\item Каков порядок группы $m\br3m$? Сколько у неё подгрупп с порядком 5?
\item Может ли пьезоэлектриком быть кристалл класса $m\br3m$? Класса $\br43m$?
\item Почему нет тетрагональной $C$-решётки Браве?
\item Чем винтовая ось $4_1$ отличается от $4_3$?
\item Какой предельной группой описывается симметрия однородного
магнитного поля и почему она не совпадает с симметрией электрического?
\end{proverka}
